%%%%%%%%%%%%%%%%%%%%%%%%%%%%%%%%%
%
%       EXERCÍCIO
%
%%%%%%%%%%%%%%%%%%%%%%%%%%%%%%%%%

\ifdefstring{\atividade}{prova}{%
    \ifdefstring{\modo}{objetivo}{%
        \renewcommand{\valorquestao}{\ValorQObj\ ponto}
    }{%
        \renewcommand{\valorquestao}{\ValorQDisc\ pontos}
    }%
}%

\begin{exercicioBanco}[\valorquestao]
Uma mostra científica deve ser exposta ao público
de acordo com as seguintes orientações:
%
\begin{enumerate}[\(\bullet\)]
    \item apresentar todos os 3 trabalhos de álgebra e todos os 3 trabalhos de geometria.
    %
    \item os trabalhos devem estar alocados em uma parede, lado a lado, de modo que os trabalhos de álgebra não fiquem de forma adjacente.
\end{enumerate}
%
De quantos modos essa mostra de trabalho pode ser exposta?

% Define as alternativas
\newcommand{\alternativas}{%
    \begin{center}
        \begin{tabularx}{\textwidth}{XXXXX}
            (a) \(6\). &
            (b) \(36\). &
            (c) \(48\). &
            (d) \(128\). &
            (e) \(144\).
        \end{tabularx}
    \end{center}
}

% Define a resposta correta
\newcommand{\resposta}{E}

% Lógica condicional para exibição
\ifdefstring{\atividade}{lista}{%
    \alternativas
    \vspace{0.5em}
    
    \noindent\textbf{Resposta:} letra \textbf{\resposta}.
}{%
    \ifdefstring{\modo}{objetivo}{%
        \alternativas
    }{%
        % modo = discursiva → não mostra alternativas
    }
}
\end{exercicioBanco}

